\subsection{Claude Code}
Anthropic's coding harness, deconstructed from its documentation as read on 15 October 2026 \citep{claudecode}; its source is not published. The loop is the open one of Algorithm~\ref{alg:harness}: the model is called until it returns a response with no tool calls, and a turn's result carries why it ended (done, a turn limit, a cost budget, an error). \prim{Input} has one entry point for everything that starts a turn: typed prompts, piped input, messages queued mid-turn and delivered after the current tool batch, scheduled prompts, and reports from background sub-agents. \prim{Context} is built as a stable cached prefix (the system prompt, the tools, then project files such as \texttt{CLAUDE.md}) with later updates appended as messages or system reminders; tool schemas for MCP servers are deferred behind a search tool; and compaction, on reaching the window, is a separate request that re-injects the project files, memory, recently read files and invoked skills. \prim{Model interface} is the Messages API with streaming, up to ten retries that never re-issue after a completed tool block, idle watchdogs, and a fallback chain the user configures. \prim{Output} is a fixed built-in tool set plus MCP tools; read-only calls run concurrently and mutating calls in sequence; shell commands run under an OS sandbox with a domain-allowlist proxy; large results spill to files; file edits are checkpointed for rewind. \prim{Control flow} adds a six-step permission pipeline (hooks, deny rules, ask rules, the permission mode, allow rules, then a prompt), delegation through an agent tool to a depth of three, and, at scale, a model-written JavaScript workflow the runtime executes, which sequences model calls and is therefore \prim{Control flow} by the rule on model-written programs.

Three components were hard to place. The web-fetch tool calls the model itself to condense a page before the main model sees it: the call is \prim{Output}, the condensing is a \prim{Context} decision, and the extra request is \prim{Model interface}. The default permission gate in ``auto'' mode is a second model asked to classify each action: a \prim{Control flow} gate implemented as a model call with its own stripped context. And a server-side variant of that classifier runs inside the API during the session's own requests, so the harness only consumes its verdict; we return to this in \S\ref{sec:summary}.

\subsection{Codex CLI}
OpenAI's terminal and IDE coding agent, at commit \texttt{f0c5077} \citep{codexcli}, written in Rust. A session is one thread that receives typed operations over a channel; a dispatcher matches each one (a turn's input, an interrupt, an approval, the answer to a question, a compaction request) and runs each user turn as a task. \prim{Input} is that channel: text typed during a running turn is queued as steering rather than starting a new turn, an inter-agent mailbox and asynchronous hook results merge into the same pending-input queue, and the model's questions to the person come back as a typed answer. \prim{Context} is the Responses API's own item list, with the environment and instructions injected as typed ``world-state'' sections sent as diffs when they change, project files walked from the root to the working directory, and tool outputs truncated in the live history only; compaction has three implementations (a local summary, a remote one performed by the server, and a token-budget reset), chosen by explicit threshold arithmetic, and can run speculatively in the background once usage passes a percentage of the window. \prim{Model interface} is the Responses API only, WebSocket first with a session-wide fallback to HTTP, five retries with jittered backoff that honour the server's retry timing, an idle timeout, and requests sent without server-side storage, with encrypted reasoning replayed. \prim{Output} routes every risky action through one orchestrator (approval, then the platform sandbox, then the run, then a retry if the sandbox denied it), with Seatbelt on macOS, seccomp with bubblewrap and Landlock on Linux and a restricted token on Windows; tools run concurrently under a read--write gate that lets parallel-safe tools overlap; and a ``code mode'' runs model-written JavaScript in V8, each nested tool call re-entering the harness's approval and sandbox path. \prim{Control flow} is one task per turn, looping until the model emits no tool call, with stop hooks, roll-over into compaction mid-turn, an approval precedence of hooks, then an optional model reviewer (``Guardian''), then the person, and sub-agents that are full sessions of their own with depth and thread caps.

Hard cases: the Guardian is a \prim{Control flow} gate whose verdict comes from a nested harness; the remote compaction is a \prim{Context} decision whose summarizer is the provider; over WebSocket the request omits the tools and instructions the server already holds, so what is sent depends on transport state while what the model is shown does not; and the plan-mode parser, which reads the streamed text into a proposed-plan item and can redirect it, is \prim{Output} rendering that also feeds the stop decision. We record the last as straddling \prim{Output} and \prim{Control flow} by design.


\subsection{opencode}
An open-source coding harness by Anomaly, at commit \texttt{654b79b} \citep{opencode}; it was rewritten since earlier versions, and the deconstruction describes the current one. Each session is an event-sourced record in SQLite, and one fibre per session drains a durable inbox: \prim{Input} is a row in that inbox, from a person, a shell command, a skill, or the harness itself, marked for delivery either at the next step boundary (``steer'') or after the current turn (``queue''). \prim{Context} re-projects the event history into model messages on every call, with instructions placed in a cached prefix and changes to them frozen into the history as dated events; compaction is triggered at about ninety percent of the window, using provider-measured token counts plus a local estimate, keeps the newest fifteen thousand tokens verbatim, and can be delegated to the provider's native compaction. \prim{Model interface} is an in-house request type compiled per provider protocol (Anthropic Messages, OpenAI Responses, Bedrock Converse, Gemini), streamed over HTTP or a stateful WebSocket, with a retry policy split by whether output had already started. \prim{Output} runs tool calls concurrently as they stream, and most tools are reachable only through one \texttt{execute} tool that runs model-written JavaScript in an interpreter that cannot call the model, which the rule on model-written programs places in \prim{Output}. \prim{Control flow} continues while a step made tool calls, stops on a text-only finish, enforces an optional per-agent step budget, and delegates to child sessions (depth one by default), whose results return as tool text or as a synthetic message.

Hard cases: every tool asserts its own permission, so an approval gate the harness consumes lives inside the tool's code (\prim{Control flow} by function); the step budget is enforced by appending an assistant message and disabling tools, a \prim{Control flow} decision carried out by editing what the model sees; and the compaction service both decides what to summarize and issues the model call itself. Native provider-side compaction leaves the harness's \prim{Context} holding an item it cannot inspect.

\subsection{pi}
A terminal coding agent by Mario Zechner, at commit \texttt{5ab7fe7} \citep{pi}, built in layers: a minimal agent loop, a provider layer, and a coding-agent session on top. \prim{Input} comes from a terminal editor, a print mode, or a JSON-RPC mode, all into one \texttt{prompt()} call, with mid-run messages split into steering (after this turn) and follow-ups (after the agent would have stopped). \prim{Context} is a projection of a persisted session tree: a compaction entry, the kept tail, and context edits, with the system prompt kept as named sections and diffed into patches, so that replaying the file reproduces the exact request; compaction triggers when the context exceeds the window less a reserve, and keeps the most recent twenty thousand tokens. \prim{Model interface} is one streaming function behind about fifteen provider adapters, with the SDKs' own retries disabled so that the harness owns a two-tier retry policy, and an overflow classifier that routes a too-long request to compaction rather than to a retry. \prim{Output} is a set of plain tool objects run in parallel by default, with mutations to one file serialized; there is no sandbox in the core (OS sandboxing and permission prompts are extensions), and a ``code mode'' tool runs model-written scripts in a QuickJS sandbox. \prim{Control flow} is a two-level loop (tool calls and steering inside, follow-ups outside) plus a session-level loop for retries, overflow recovery and extension continuations; sub-agents are not in the core, though an example spawns child processes; and per-request routing between ``virtual models'' is keyed to the reason for the request.

Hard cases: a wire-level hook lets an extension rewrite the request body, which is \prim{Context} by function though it sits in the adapter; the virtual-model router has a ``retry'' branch that overlaps the backup-model rule (we place it in \prim{Control flow}, since the router decides by the state of the work); and the code-mode script can call auxiliary models from inside the sandbox, a model call issued from inside a tool, which we discuss in \S\ref{sec:summary}. A telemetry interface exists, but no span is emitted at this commit.

\subsection{Hermes Agent}
A personal agent by Nous Research, at commit \texttt{7bf3766} \citep{hermes}, with one agent core behind a terminal interface, about twenty-five messaging-platform adapters, an HTTP server, an ACP adapter and a cron scheduler. \prim{Input} is one normalized message event shared by all of these; mid-turn messages become steering rather than new turns. \prim{Context} is a three-tier system prompt, kept byte-stable so that caching holds: identity, then project files (capped by size), then volatile memory, skills and the time; compaction at half the window prunes, then summarizes with an auxiliary model; oversized tool output is spilled to files rather than dropped; and a tool-search bridge lets the model fetch tool schemas on demand. \prim{Model interface} is an OpenAI-style chat-completions client by default with pluggable transports for Anthropic, Responses and Bedrock, streaming on, three retries with backoff, then a configurable fallback chain, then a bounded auto-recovery wait. \prim{Output} sends a fixed core tool set on every call, splits each batch into parallel-safe segments, runs terminal commands behind pluggable sandboxes (local, Docker, Modal, and others) and a regex-plus-human approval gate, and persists each tool-call row before executing it. \prim{Control flow} is a loop of named phases, each returning a verdict, with an unlimited iteration budget by default, ``stop-gate'' nudges that append a synthetic message and go again, flat delegation to threaded children, and task-keyed routing to auxiliary models.

Two components do not fit cleanly. In one mode the whole turn is handed to an external Codex process, and Hermes keeps only input, persistence and a fallback: the five primitives collapse into one opaque external harness that is neither a sub-agent (there is no outer loop) nor the model. And ``providers'' reached over ACP make the thing behind the \prim{Model interface} another agent. Both are nested harnesses by the rule of \S\ref{sec:rules}, but of an unusual shape: the nesting is at the \prim{Model interface}, not in \prim{Control flow}.

\subsection{A LangGraph workflow}
LangGraph \citep{langgraph}, at commit \texttt{5eb60ae}, is a framework for writing control flow as a graph over a checkpointed state, executed by a runtime in bulk-synchronous supersteps. We deconstruct the agent it ships, \texttt{create\_react\_agent}, together with that runtime. \prim{Input} is a write to a named channel of a typed state, merged by a user-chosen reducer, whether it is a person's message, a tool result, or a resume value after an interrupt. \prim{Context} is exactly the state's message list, optionally prefixed by one system message; fitting it to the window is left entirely to a user-written hook before the model call, and the framework provides only a primitive for removing messages. \prim{Model interface} is one call to a LangChain chat model per step, and the framework contains no client of its own, no streaming of its own, and no backup model; retries and timeouts are per-node policies, off by default. \prim{Output} is a tool node that runs every tool call of a response concurrently in-process, turns results and (configurably) errors into messages, and lets a tool return a routing command. \prim{Control flow} is the graph: ``go again'' is a conditional edge from the agent node to the tool node and back, ``stop'' is an empty task set or a recursion limit, and a pause is an exception the runtime persists as a checkpoint write and resumes by re-running the node.

This is the clearest case of the three-level distinction. The framework fixes the five primitives' mechanisms and leaves most of the design decisions to user code: the edges, the stopping rule, the window-fitting hook, the approval gate around an interrupt, dynamic model selection, retry policies. A workflow and an agent built with it differ in the graph, which is \prim{Control flow}, and nowhere else. The hard cases are the step counter injected into the state (it looks like \prim{Context} but no model sees it; it is a \prim{Control flow} budget), and a scheduling packet that carries a state snapshot to the tool task (\prim{Control flow} carrying \prim{Context}).
